\scanpage{166}
\mspage{82}
% archive sheet number: 165
% margin note written slanted
\margin{Introduire $I(J)$ \unsure{avant} l'énoncé de la proposition. \uline{NB} La condition de cofinalité $I(J)\longrightarrow J$ n'est pas automatique, comme on voit \unsure{par l'exemple trivial} $J=\widehat{B}_G$, $G\ne1$.}
\margin{\uline{NB} L'hyp. de cofinalité superflue si $J$ est rigide, i.e. les \unsure{endom.} d'un objet de $J$ sont identiques.}

% lines 1-5, continuation of condition d)
ordonnons $I$ par inclusion. Alors $I$ est grand devant $\pi$, de plus \deleted{\unsure{et}} \deleted{$J$ est filtrante. (Cette dernière condition est superflue si $J$ est rigide, i.e. \unsure{les} \unsure{endom.} d'un \unsure{objet de $J$} sont l'identité).} \inserted{et le foncteur canonique $I\longrightarrow J$ ($J'\longmapsto e_{J'}=$ objet final de $J'$) cofinal.}

\uline{Démonstration}

% lines 6-9
$a\Longrightarrow b$ vérification immédiate, \deleted{\unsure{en appliquant}} \inserted{en utilisant} \deleted{\unsure{de la}} condition F1, F2 de cofinalité pour \deleted{des catégories} $I$ filtrantes

$b\Longrightarrow c$ par AQT

$c\Longrightarrow d$ par AQT

% lines 12-17
$d\Longrightarrow a$ Considérons l'ens. ordonné grand devant $\pi$ construit dans l'énoncé de d), et le \uline{foncteur}
\[
I\longrightarrow J
\]
donné par
\[
J'\longmapsto e_{J'}=\text{objet final de }J'.
\]
(Vérifier que c'est un foncteur: \unsure{immédiate|immédiat}.) Ce foncteur est surjectif, donc satisfait la condition F1 de cofinalité.

% lines 18-24
La condition F2 résulte aussi de d)
\diagram{\begin{tikzcd}[ampersand replacement=\&]
j \arrow[r,shift left=1.3ex,"u_1"] \arrow[r,shift right=1.3ex,"u_2"'] \& e_{J'}
\end{tikzcd}}
$J'$, en regardant les trois objets $J'$, $j\xrightarrow{u_1}e_{J'}$, $j\xrightarrow{u_2}e_{J'}$ de $I$, et en prenant un $J''$ qui les majore tous trois, alors $e_{J'}\longrightarrow e_{J''}$ égalise $u_1,u_2$.
% margin note written slanted
\margin{Un peu bref le cas où $e_{J'}=j$ \unsure{échappe} à cet argument, et \unsure{demande} plus.}

% line 25, wholly cancelled
\deleted{\uline{Corollaire} Soit $I\longrightarrow J$ un foncteur cofinal.}

% lines 26-29
\uline{Corollaire} 4.4.5. Soit $F\in\operatorname{Ob}\widehat C$. Pour que $F$ appartienne à $\operatorname{Ind}'_{\pi}(C)$, il f. et s. que $C_{/F}$ satisfasse à l'une des quatre conditions équivalentes de 4.4.4.
% vertical marginal emphasis stroke
\margin{$\big|$}

% lines 30-33
\uline{Cor.} 4.4.6. Soit $J$ une catégorie stable par $\unsure{\underline{\lim}|\varinjlim}$ de \deleted{\unsure{type}} cardinal $\leq\pi$. Alors $J$ est grande devant $\pi$. En effet, les conditions $\mathrm{PF}_{\pi}1$, $\mathrm{PF}_{\pi}2$ sont \unsure{visiblement} vérifiées.
